\chapter{相对论中的能量与动量}

\section{相对论与哲学家}

在这一章中，我们将继续讨论爱因斯坦和庞加莱的相对性原理，因为它们影响着我们的物理观念以及人类思维的其他分支。

庞加莱以如下方式表述了相对性原理：“按照相对性原理，对于一个固定的观察者与对于一个相对于他作匀速运动的观察者来说，描述物理现象的定律必须是相同的，因而我们没有，也不可能有任何一种方法去辨认我们是否参与了这样一种运动。”

当这种观念披露于世时，在哲学家中引起了很大的骚动，特别是那些“鸡尾酒会哲学家”。他们说：“噢，这很简单：爱因斯坦的理论表明，一切都是相对的！”事实上，不仅是在鸡尾酒会上所见到的那些哲学家（为了不使他们难堪，我们就称他们为“鸡尾酒会哲学家”），而且数目多得令人吃惊的哲学家们都纷纷声称：“一切皆相对，此乃爱因斯坦之推论，它给子吾等之观念以深远的影响。”他们还补充说：“物理学亦已表明，现象有赖于人们的参照系。”诸如此类的话我们已听得很多，但是要弄清楚它们的含义则非常困难。大概原来所指的参照系，就是指我们在对相对论的分析中所用到的坐标系。这样，“事物有赖于人们的参照系”这个事实，就被设想为曾给于现代观念以深刻的影响。人们很可能对此感到不解，因为归根结底，事物依赖于一个人所抱的观点这件事是如此简单，为了要发现它，肯定不会有什么必要到物理学的相对论中去找麻烦。任何一个在街上散步的人肯定都明白，他所看到的一切取决于他的参照系，因为当一个过路人走近他时，他首先看到的是那个人的前面，而后再看到其后面；在据说是源出于相对论的大多数哲学中，没有比“一个人从前面看与从后面看不同”这种说法更深刻的了。几位盲人把大象描写成几种不同的样子，这个古老的故事或许是哲学家对相对论所抱有的观点的另一个例子

但是在相对论中，一定具有比刚才“一个人从前面看与从后面看不同”那种简单的说法更为深刻的含义。相对论当然要比这种说法深刻得多，因为我们能够借助于它作出确定的预言。如果单从这样简单的观察居然能够预知自然界的行为，那么它一定是相当令人吃惊的。

也有另一学派的哲学家，他们对于相对论感到很不舒服，因为相对论断定如果不往外看，我们就无法确定我们运动的绝对速度，而这些人则说：“一个人不往外看就不能测出他的速度，那是很明显的。不往外看面去谈论一个物体的速度毫无意义，这也是不证自明的一件事：物理学家相当笨抽，因为他们的想法不是这样，可是现在总算使他们明白过来情况就是这样。只要我们哲学家认识到物理学家所思考的问题是什么，那么我们就会立即通过大脑来判断，不往外看就不可能说出一个人的运动有多快，这样我们就会对物理学作出巨大贡献了。”这样的哲学家总是有的，他们在我们周围喋喋不休地企图告诉我们一点什么东西，但是，实际上他们从未理解过这类问题的细致和深刻之处。

我们之所以无法鉴别绝对运动，乃是实验的一个结果，而不是像我们所能很容易想象的那种只是单纯思维的结果。首先，牛顿就已确实认为，假如一个人沿直线作匀速运动，那么他就不能说出自己跑得多快。事实上，牛顿最早说明了相对性原理，上一章中的一段引文就是他对此的陈述。那么，为什么在牛顿时代哲学家们就没有提出“一切都是相对的”或其他别的什么来喧哗一番呢？这是因为直到麦克斯韦提出电动力学理论，才有物理定律认为人们不往外看能够测定他的速度；不久就从实验上发现，这是不行的。

那么，一个人不往外看就无法知道他究竟运动得多快这一点是绝对的，肯定的，哲学上必然的吗？相对论的结果之一是发展了一种哲学，这种哲学认为：“你只能定义你所能测量的东西！因为非常明显，一个人如果不去看他相对于什么在测量速度，那么就无法测量这个速度，所以很清楚，谈论绝对速度是毫无意义的。物理学家应该领会到，他们所能谈论的只是那些他们所能测量的东西。”但是整个问题在于：人们是否能定义绝对速度这个问题是与下一个问题相同的，即人们是否在一个实验中能不往外看就觉察到他是否在运动。换句话说，某一事物是否可以测量，并不是由纯粹思维所能先验地予以决定，而是只能由实验来决定。假定我们已知光的速度为$186000~mi \cdot s^{-1}$，那么人们将会发现，没有几个哲学家会沉着地说：“这是不证自明的，如果光在汽车中的传播速度为$186000~mi \cdot s^{-1}$，汽车的车速为$100000~mi \cdot s^{-1}$，那么对于地面上的观察者来说，光的速度也是$186000~mi \cdot s^{-1}$。”对于这些哲学家来说这是一个令人吃惊的事实：正是这些认为这是很明显的人，当你告诉他们一件特殊的事实的时候，他们就认为这不是那么明显了。

最后，甚至还有这么一种哲学，它认为除非我们往外看，否则我们就不能觉察任何运动。在物理学中这是根本不正确的。诚然，人们无法觉察沿直线的匀速运动，但是如果整个房间在转动，那么我们就一定能知道它，因为每个人将被掷到墙上——这里有各种各样的“离心”效应。利用所谓傅科（Foucsult）摆的方法，地球的绕轴转动可以不用观察星体而加以确定。因此，“一切都是相对的”这句话并不正确；只有匀速运动在不往外看时是觉察不到的。绕固定轴的匀速转动是可以觉察到的。当把这一点告诉一个哲学家时，他会十分心烦意乱：他实在不能理解这件事，因为在他看来，不往外看而能确定绕轴转功似乎是不可能的。如果这个哲学家是一个相当出色的哲学家，那么过一段时间后他可能会回过头来对我们说：“我明白了！我们确实并没有绝对转动这样一种事：你知道，实际上我们是相对于星体在转动。因而星体对物体所施加的某种影响必然引起了离心力。”

现在就我们所知的一切情况来说，那是对的；目前我们还没有办法可以告诉你，如果周围没有星体成星云，那么是否还会存在离心力。我们没有可能去做这样一个实验，先把所有的星云移开，而后去测量地球的转动，所以我们也就不知道。我们必须承认这位哲学家可能是对的。因此他高兴地回转身来说：“世界最终变成这个样子是绝对必要的，即绝对的转动是毫无意义的，它只是相对于星云而言的。”于是我们对他说：“那么，我的朋友，相对于星云所作的匀速直线运动，不应该在一辆汽车内产生任何效应这一点是明显的还是不明显的呢？”现在，运动已不再是绝对的，而是相对于星云的，它变成了一个神秘的问题，即一个只能用实验来回答的问题。

那么，什么是相对论对哲学的影响呢？如果我们只限于去谈这种意义上的影响，例如相对论原理给物理学家带来了哪些新的观念和启示，那么我们可讲一下其中的如下几点。首先，我们的一个发现主要是，有些概念即使在长时期内都被认为适用，并且得到了十分准确的验证，然而它们仍然可能是错误的。牛顿定律在这么多年被视为似乎是正确的之后，居然被说成是错误的，这显然是一件使人震惊的发现。当然，很清楚不是实验错了，而只是因为这些实验是在有限的速度范围内完成的，而这些速度之小，不可能使相对论效应明显地表现出来。但无论如何，我们现在对于物理定律已抱有一种远为谦逊的见解，即任何一件事都可能是错的！

第二，如果我们有一些“奇特”的观念，比如时间会随着运动而变慢，等等，那么，究竟我们是喜欢它们还是不喜欢它们，则是与此不相干的另一个问题。唯一与之有关的问题，便是这些观念是否与实验上的发现相一致。换句话说，“奇特的观念”只要符合于实验就行，而我们必须讨论钟的行为等等，其唯一的理由就是要说明，虽然时间膨胀的概念是多么奇怪，但它与我们测量时间的方式是协调的。

最后，第三个启示虽然略带一些技术性，但在我们研究其他物理定律时它证明是非常有用的，这就是要注意定律的对称性，或者更明确地说，就是要寻找这样一种方式使得定律在变换时能保持其原有形式不变。在讨论矢量理论时，我们曾看到，在转动坐标系时，运动的基本定律并没有改变，而现在我们又知道，在以一种特有的由洛伦兹变换所提供的方法改变时间和空间变量时，它们的形式也没有改变。因而，这个观念，即在什么形式或作用下基本定律仍保持不变，已被证明是一个非常有用的观念。

\section{孪生子佯谬}

为了继续讨论洛伦兹变换和相对论效应，我们来考虑一个著名的所谓彼得（Peter）和保罗（Paul）的“伴谬”，并假定他们是同时出生的一对李生子。当他们成长到能操纵宇宙飞船时，保罗以很快的速度飞了出去，彼得则仍留在地面。由于彼得看到保罗运动得这么快，所以从彼得的观点看来，保罗的钟似乎走慢了，他的心跳变慢了，思维也迟缓了，每件事都延迟了。当然，保罗自己并没有感到出现任何异常情况，但是如果他在外面漫游了一段时间之后再回到地面，他将比在地面上的彼得年轻！这确实是对的，它是相对论的结论之一，而相对论是被清楚地证实了的。正像子运动时，它的寿命要延长一样，当保罗运动时，他的寿命也会延长。这件事只有在这些人眼光中才称为是一种“谬”，因为他们认为，相对论原理意味着一切运动都是相对的，他们说：“好，好，好，从保罗的观点看来，难道我们不是也可以说彼得正在运动，因而他应当衰老得慢一点吗？由于对称性，唯一可能的结果是，当他们会面时，大家的年龄应当相同。”但是，为了使他们能重新相遇并进行比较，保罗必须要么在旅途的终点停下来，并且将钟进行比较，要么更简单一些，他必须返回，而返回的那个人必定是正在飞行（或运动）的那个人，他知道这一点，因为他必须转过身来飞行。当他转过身来的时候，他的飞船上各种不寻常的事情就发生了——火箭射了出去，东西向墙上撞了过去，等等——而彼得则一点也没有感到什么。

所以，如果要叙述这条规则的话，就可以说：感觉到加速度和看到东西向墙上撞了过去等等的那个人，将是比较年轻的一个；这就是他们之间在“绝对”意义上的一个差别，而这肯定是正确的。当我们讨论运动的$\mu$子的寿命变长这个事实时，作为例子我们使用了它们在大气中的直线运动。但是我们也可以在实验室里产生$\mu$子，并用磁铁来使它们作曲线运动，即使在这种加速运动下，它们的寿命的延长与在直线运动的情况下完全一样。虽然还没有一个人具体地安排过一个实验，使我们能消除这个佯谬，但是我们可以把一个静止的$\mu$子同一个跑完整个一圈的$\mu$子相比较，这时肯定将会发现绕过整个一圈的$\mu$子的寿命要长一些。虽然我们实际上还没有用整个一圈做过这样的实验，但其实这并不必要，因为一切事情都符合得很好。对于那些坚持认为每个单独的事实都要直接得到证实的人，这或许不能使他得到满足。但是我们有充分把握来预言保罗转过整个一圈的那个实验的结果。

\section{速度的变换}

爱因斯坦相对性与牛顿相对性之间的主要差异，在于把两个处于相对运动中的系统之间的坐标与时间联结起来的变换规律是不同的。正确的变换规律，也就是洛伦兹变换方程为
\begin{equation}
    \label{Eq:I:16:1}
    \begin{aligned}
        x'&=\frac{x-ut}{\sqrt{1-\frac{u^2}{c^2}}} \\
        y'&=y \\
        z'&=z \\
        t'&=\frac{t-\frac{ux}{c^2}}{\sqrt{1-\frac{u^2}{c^2}}}
    \end{aligned}
\end{equation}
这些等式对应于一种比较简单的情况，即其中两个观察者的相对运动沿着它们共同的$x$轴。当然也有可能沿着其他方向运动，但是最一般的洛伦兹变换则将由于所有四个量都混杂在一起而变得相当复杂。我们将继续使用这个较为简单的形式，因为它包含了相对论的所有主要特点。

我们现在来进一步讨论这个变换的推论。首先，有趣的是倒过来解这些等式。这就是说，这里有一组线性方程，它们是四个方程和四个未知数，可以倒过来解这些方程，即用$x'$，$y'$，$z'$，$t'$来表示$x$，$y$，$z$，$t$。结果是非常有意思的，因为它告诉我们，从“运动”坐标系的观点来看，“静止”坐标系会是什么样子。当然，由于运动是相对的和匀速的，所以“运动”的那个人，如果愿意的话也可以说：实际上是另一个人在运动，而他自己则静止着。因为这另一个人是在沿着反方向运动，所以他应当得到同样的变换，但是速度要用相反的符号。这与我们在演算中所得到的完全相同，因而是协调的。假如得出的结果不是如此，那我们倒真有理由要担心了！
\begin{equation}
    \label{Eq:I:16:2}
    \begin{aligned}
        x&=\frac{x'+ut'}{\sqrt{1-\frac{u^2}{c^2}}} \\
        y&=y' \\
        z&=z' \\
        t&=\frac{t'+\frac{ux'}{c^2}}{\sqrt{1-\frac{u^2}{c^2}}}
    \end{aligned}
\end{equation}

其次，我们来讨论相对论中速度的叠加这个有趣的问题。我们还记得一个奇特的难题，即对于所有系统来说，光速都是$186000~mi \cdot s^{-1}$，即使它们处于相对运动之中也是如此。像如下例子所说明的那样，这是较普遍的问题中的一个特殊情况。假定宇宙飞船内有一个物体以$10000~mi \cdot s^{-1}$运动，飞船本身的速度是$10000~mi \cdot s^{-1}$；从外部观察者的观点看来，飞船内这个物体是以多大的速度在运动？我们也许要说$200000~mi \cdot s^{-1}$，这就要比光速还快。这是非常使人沮丧的，因为不能设想它会跑得比光速还快！普遍的问题则如下所述。

假设在飞船内有一个物体，从飞船内一个人的观点看来，它以速度$v$运动，而飞船本身则相对于地面有一速度$u$。我们所要知道的是从地面上一个观察者的观点看来，这个物体以多大的速度$v_x$运动。当然，这还是一个在$x$方向上运动的特殊情况。此外，还存在着在$y$方向或任何方向上的速度变换；这些都可以在需要时加以导出。在飞船内物体的速度是$v_{x'}$这就说明位移$x'$等于速度乘以时间，即
\begin{equation}
    \label{Eq:I:16:3}
    x'=v_{x'}t'
\end{equation}
现在我们只要对一个在$x'$和$t'$之间具有关系式（\ref{Eq:I:16:3}）的物体去计算它从外部观察者的观点看来的位置与速度是多少，因而我们只要简单地把式（\ref{Eq:I:16:3}）代入式（\ref{Eq:I:16:2}），就得到
\begin{equation}
    \label{Eq:I:16:4}
    x=\frac{v_{x'}t'+ut'}{\sqrt{1-\frac{u^2}{c^2}}}
\end{equation}
但是我们在这里看到$x$是用$t'$表示的。为了得到外部观察者所看到的速度，我们必须用他的时间，而不是用另一个人的时间去除他的距离！所以我们也必须计算外部观察者所看到的时间，即
\begin{equation}
    \label{Eq:I:16:5}
    t=\frac{t'+\frac{u(v_{x'}t')}{c^2}}{\sqrt{1-\frac{u^2}{c^2}}}
\end{equation}
现在必须找出$x$与$t$之比，这就是
\begin{equation}
    \label{Eq:I:16:6}
    t=\frac{x}{t}=\frac{u+v_{x'}}{1+\frac{uv_{x'}}{c^2}}
\end{equation}
其中平方根已被消去。这就是我们所要求的定律：合速度，即两个速度之“和”，并不恰巧是两个速度的代数和（我们知道，不可能是两者的代数和，否则就会使我们陷于困境），而是为（$1+\frac{uv}{c^2}$）所“校正”了的。

现在我们来看看将会发生什么。假定你们在宇宙飞船内以光速的一半在运动，而飞船本身也在以光速的一半飞行。因此，$u=\frac{c}{2}$，$v=\frac{c}{2}$，而分母中的$\frac{uv}{c^2}=\frac{1}{4}$，于是
\begin{equation*}
    v=\frac{\frac{1}{2}c+\frac{1}{2}c}{1+\frac{1}{4}}=\frac{4}{5}c
\end{equation*}
所以在相对论中，“$\frac{1}{2}$”加“$\frac{1}{2}$”并不等于1，而只等于“$\frac{4}{5}$”。当然，低速可以很容易用熟悉的方法相加，因为只要速度与光速相比很小，我们就可忘掉（$1+\frac{uv}{c^2}$）这个因子；但是在高速情况下，事情就完全不同，而且也变得非常有趣了。

我们来考虑一个极限情况。开一个玩笑吧！宇宙飞船里的那个人正在观察光本身。也就是说$v=c$，而飞船是在以$u$运动。那么从地面上的人看来将会怎样？回答是
\begin{equation*}
    v=\frac{u+c}{1+\frac{uc}{c^2}}=\frac{c(u+c)}{u+c}=c
\end{equation*}
因此，如果飞船里有什么东西在以光速运动，那么从地面上的观察者的观点看来，它还是以光速运动！这很好，因为事实上，这正是爱因斯坦相对论首先打算要做到的——可见这个理论不错！

当然，也有一些情况，运动并不是在匀速移动的那个方向上进行的。比如飞船中可能有个物体正好相对于飞船以速度$v_{y'}$“朝上”运动，而飞船则在“水平”地飞行着。这时，我们只要按照同样的方法去做，不过用$y$项而不用$x$项就是了，结果是
\begin{equation*}
    y=y'=v_{y'}t'
\end{equation*}
所以，如果$v_{x'}=0$，则
\begin{equation}
    \label{Eq:I:16:7}
    v_y=\frac{y}{t}=v_{y'}\sqrt{1-\frac{u^2}{c^2}}
\end{equation}
因此，侧向的速度不再是$v_{y'}$，而是$v_{y'}\sqrt{1-\frac{u^2}{c^2}}$。这个结果我们是用代入与组合变换方程得到的，但我们也能由于如下理由而直接从相对论原理看出这个结果（再去探索一下总是好的，看看我们是否能找出其理由来）。我们已经讨论过（图\ref{figure:光束经过的斜向路程示意图}）一只钟在运动时如何走动
\begin{wrapfigure}{r}{0.4\textwidth}
    \centering
    \includegraphics[width=0.38\textwidth]{Chapter16/光和“粒子”在一个运动的钟内的运动轨迹}
    \caption{光和“粒子”在一个运动的钟内的运动轨迹}
    \label{figure:光和“粒子”在一个运动的钟内的运动轨迹}
\end{wrapfigure}
的；在固定的地面坐标系看来，光线以速度$c$沿一定偏角前进，而在运动坐标系看来，则只是以同样速度沿垂直方向运动。我们曾经看到，在固定坐标系中光速的垂直分量比光速本身小一个因子$\sqrt{1—\frac{u^2}{c^2}}$（见图\ref{figure:光束经过的斜向路程示意图}）。但现在我们假定，让一个实物粒子在这同一只“钟”内来回跑动，其速度为光速的整分数$\frac{1}{n}$（图\ref{figure:光和“粒子”在一个运动的钟内的运动轨迹}），那么当粒子来回跑一次时，光将恰好走过$n$次。也就是说：“粒子”钟每一次的“滴答”声恰好与光钟的第$n$次“滴答”声相符合。当整个系统在运动时，这个事实仍然正确，因为符合一致的物理现象在任何参照系中仍将是符合一致的。因此，由于$c_{y}$小于光速$c$，粒子的速度$v_y$也必然要比对应的速度小同一个平方根因子！这就是平方根所以会出现在任何一个垂直速度中的原因所在。

\section{相对论性质量}

在上一章中我们已经看到一个物体的质量随着其速度的变大而增加，但没有对此加以说明，这就是说，我们没有进行过类似于对钟的行为所作出的那样的论证。然而，我们能够证明，作为相对论加上少数几个其他合理的假设的结果，质量是必须按照这种方式变化的我们必须说“少数几个其他假设”，因为只要我们还打算进行有意义的推理，那么，除非假定有某些定律已经成立，不然就不可能证明任何东西）。为了避免要去研究力的变换定律，我们将研究碰撞这个问题，在那里，除了假设动量和能量都守恒外，不需要知道有关力的任何定律。同样，我们将假设运动粒子的动量是一个矢量，而且总是指向速度的方向。然而我们将不像牛顿所做的那样，把动量设为一个常熟乘上速度，而只是把它设为速度的某一函数。因此我们把动量矢量写成某一个系数乘上速度矢量
\begin{equation}
    \label{Eq:I:16:8}
    \boldsymbol{p} = m_v\boldsymbol{v}
\end{equation}
我们在系数上记一个下标“$v$”，为的是要提醒大家，它是速度的一个函数，而且我们也同意把这个系数$m$，称为“质量”。当然，当速度很小时，它与我们过去在低速实验中所测得的质量是相同的。现在，我们试图从物理定律必须在每个坐标系中都相同这个相对论原理来论证$m$。的公式必须是$\frac{m_0}{\sqrt{1-v^2/c^2}}$。

假定我们有两个粒子，比如两个质子，它们完全相同，并且以精确相等的速率相向运动，它们的总动量为零。现在，可能出现什么情况呢？碰撞以后，它们的运动方向必须正好相反，因为如果不是正好相反，那么总的动量矢量就不会是零，动量也就不会守恒。它们还必须具有相同的速率，因为彼此之间是完全相类似的；事实上，它们的速率必须都与碰撞前相同，因为我们假定能量在这些碰撞中是守恒的。所以这种弹性碰撞是一个可逆碰撞，如图\ref{figure:两个以相等速率、相向运动的相同粒子所发生的弹性碰撞的两种视图}(a)所示：所有的箭矢长度相等，所有的速率大小相等。我们假定，这种碰撞总是可以任意安排的，即在这样一种碰撞中，可以出现任何角度$\theta$，也可以用任何大小的速度。其次，我们注意到这同一个碰撞可以通过坐标轴的转动从不同的角度来观察。正是为了方便起见，我们将这样来转动坐标轴，使水平轴把它平均分成两半，如图\ref{figure:两个以相等速率、相向运动的相同粒子所发生的弹性碰撞的两种视图}(b)所示。图中只是把坐标轴转动后的同一个碰撞重新画出而已。

\begin{figure}[htbp]
    \centering
    \begin{minipage}[t]{0.4\textwidth}
        \centering
        \includegraphics[width=5cm]{Chapter16/两个以相等速率、相向运动的相同粒子所发生的弹性碰撞的两种视图}
        \caption{两个以相等速率、相向运动的相同粒子所发生的弹性碰撞的两种视图}
        \label{figure:两个以相等速率、相向运动的相同粒子所发生的弹性碰撞的两种视图}
    \end{minipage}
    \begin{minipage}[t]{0.4\textwidth}
        \centering
        \includegraphics[width=5cm]{Chapter16/从运动汽车上看上去的碰撞的另外两种视图}    
        \caption{从运动汽车上看上去的碰撞的另外两种视图}
        \label{figure:从运动汽车上看上去的碰撞的另外两种视图}
    \end{minipage}
\end{figure}

真正的技巧在于：我们从某个驱车前进的人的角度来观察这个碰撞，汽车的速度等于粒子1的速度的水平分量。那么，这个碰撞看上去像什么呢？就粒子1而言，它看上去在径直朝上跑，因为它已经没有那个水平分量，然后它又垂直往下落，这也是因为它没有那个水平分量。也就是说，碰撞看上去就如图\ref{figure:从运动汽车上看上去的碰撞的另外两种视图}(a)所示那样。然而，粒子2却按另一种方式飞行，当我们驱车经过时，它看来以更大的速度和较小的角度飞行，但是我们可以判断出碰撞前后的角度是相同的。我们以$u$表示粒子2的速度的水平分量，以$w$表示粒子1的垂直速度。

现在，问题是粒子2的垂直速度$utan\alpha$是什么？假如我们知道的话，就可以利用垂直方向的动量守恒定律来得到动量的正确表示式。很清楚，水平方向的动量分量是守恒的：对于两个粒子来说碰撞前后都相同，对于粒子1来说，水平分量为零。所以我们只需要对垂直向上的分速度$utan\alpha$应用守恒定律。但是，只要以另一种方式观察同样的碰撞，我们就可以求得朝上的速度！如果我们从以速率$u$向左边运动的车来观察图\ref{figure:从运动汽车上看上去的碰撞的另外两种视图}(a)的碰撞，那么所看到的只是将图\ref{figure:从运动汽车上看上去的碰撞的另外两种视图}(a)的情况“翻转过来”而已，如图\ref{figure:从运动汽车上看上去的碰撞的另外两种视图}(b)所示。现在粒子2以速度$w$飞下又飞上，而粒子1得到了水平速度$u$。当然，现在我们知道速度$utan\alpha$等于什么了；它就是$w\sqrt{1-\frac{u^2}{c^2}}$[参见式（\ref{Eq:I:16:7}）]。我们还知道图\ref{figure:从运动汽车上看上去的碰撞的另外两种视图}(b)中垂直运动粒子的垂直动量变化是
\begin{equation*}
    \Delta p = 2m_w w
\end{equation*}
（上式乘以2是因为它的运动先朝下再朝上。）而斜向运动的粒子具有一定的速度$v$，它的分量我们发现是$u$和$w\sqrt{1-\frac{u^2}{c^2}}$，它的质量为$m_v$，。因此这个粒子垂直动量的变化是$\Delta p'=2m_vw\sqrt{1-\frac{u^2}{c^2}}$，因为，按照我们假设的定律式（\ref{Eq:I:16:8}），动量分量等于与速度数值相应的质量乘以速度在该方向上的分量。于是为了使总动量为零，垂直方向上的两个动量必须相抵消，因此，以速度$v$运动的质量与以速度$w$运动的质量之比必须为
\begin{equation}
    \label{Eq:I:16:9}
    \frac{m_w}{m_v}=\sqrt{1-\frac{u^2}{c^2}}
\end{equation}

我们取一个$w$是无限小的极限情况。如果$w$确实很小，那么很清楚，$v$与$u$实际上是相等的。在这种情况下，$m_w \to m_0$，而$m_v \to m_u$，于是，就得到一个重大的结果
\begin{equation}
    \label{Eq:I:16:10}
    m_u=\frac{m_0}{\sqrt{1-\frac{u^2}{c^2}}}
\end{equation}
作为一个有趣的练习，我们现在检验一下，假定式（\ref{Eq:I:16:10}）是正确的质量公式，那么式（\ref{Eq:I:16:9}）是否对任意的$w$值都确实成立。注意，式（\ref{Eq:I:16:9}）中所需的速度$v$可由直角三角形算出
\begin{equation*}
    v^2=u^2+w^2(1-\frac{u^2}{c^2})
\end{equation*}
经过适当计算后，我们发现式（\ref{Eq:I:16:9}）确实成立，虽然起先我们只是在$w$很小的极限情况下利用了这个等式。

现在，我们承认动量是守恒的，质量与速度的关系是由式（\ref{Eq:I:16:10}）决定的，然后再继续看看我们是否还能得出别的什么结论。我们来考虑一种通常称为非弹性碰撞的过程。为了
\begin{wrapfigure}{r}{0.5\textwidth}
    \centering
    \includegraphics[width=0.48\textwidth]{Chapter16/质量相同的物体之间的非弹性碰撞的两种看法}
    \caption{质量相同的物体之间的非弹性碰撞的两种看法}
    \label{figure:质量相同的物体之间的非弹性碰撞的两种看法}
\end{wrapfigure}
简单起见，假设属于同一类型的两个物体，以相同的速率$w$相向运动，彼此碰撞后结合在一起成为新的静止的物体，如图\ref{figure:质量相同的物体之间的非弹性碰撞的两种看法}(a)所示。我们知道对应于$w$的每个物体的质量$m$是$\frac{m_0}{\sqrt{1-w^2/c^2}}$。如果我们假定动量守恒和相对论原理，就可以论证一件有关这个所形成的物体质量的有趣事实。我们设想有一个垂直于$w$的无限小的速度（对于有限的值也可以同样处理，但是对无限小速度的情形更易于理解），并在一个速度为（$-u$）的电梯上来观察这一碰撞。我们所见到的情况如图\ref{figure:质量相同的物体之间的非弹性碰撞的两种看法}(b)所示。复合物体的质量$M$是未知的。现在物体1以朝上的分速度$u$和实际上就等于$w$的水平分速度运动，物体2也是如此。碰撞后，质量为$M$的物体以速度$u$朝上运动，$u$与光速相比是极小的，与$w$相比也很小。由于动量必须守恒，所以我们估算一下在朝上的方向上碰撞前后的动量。在碰撞前，动量$p \approx 2m_wu$，碰撞后的动量显然为$p' = M_uu$，但由于$u$是这样小，所以$M_u$基本上可以认为是$M_0$。这些动量按照守恒定律必须相等，所以
\begin{equation}
    \label{Eq:I:16:11}
    M_0=2m_w
\end{equation}
两个相等的物体碰撞后所形成的那个物体的质量必定为进行碰撞的物体质量的两倍。你们也许会说：“当然，这就是质量守恒嘛。”但是并不那样容易，因为这些质量虽然比起它们静止时的质量来增加了，但对总的$M$，它们提供的不是处于静止时的质量，而是更多一些。令人惊奇的是，为了使两个物体碰撞时动量守恒成立，即使在碰撞后形成的物体处于静止状态，其质量也必须大于物体的静止质量！

\section{相对论性能量}

在上一章中，我们论证了作为物体质量对速度的依赖关系与牛顿定律的结果，力对一个物体所做的总功引起的动能改变总是
\begin{equation}
    \label{Eq:I:16:12}
    \Delta T = (m_u-m_0)c^2=\frac{m_0c^2}{\sqrt{1-\frac{u^2}{c^2}}}-m_0c^2
\end{equation}
我们甚至还进一步推测全部能量是总质量乘以$c^2$。现在我们来继续这个讨论。

假定在$M$内仍可以“看见”发生碰撞的两个等质量的物体。例如，一个质子和一个中子“粘合在一起”，但仍在$M$内来回运动。虽然起初预料$M$的质量是$2m_0$，但结果我们发现它并不是$2m_0$，而是$2m_w$。由于$2m_w$是形成$M$时加入的质量，而$2m_0$是在$M$中的物体的静止质量，因此复合物的过剩质量就等于带进去的动能。当然，这意味着能量有惯性。前面一章中我们讨论过气体的加热，说明了由于气体分子在运动，而运动的物体质量增大，当我们对气体加入能量后，它的分子运动就加快，因而气体质量也变大。其实上述论证完全是一般性的，我们关于非弹性碰撞的讨论表明无论动能是否存在，这部分质量总是存在的。换句话说：如果两个粒子彼此靠近，产生势能或任何其他形式的能量，或者这两部分由于越过势垒、克服内力做功等等而使速度变慢，这时质量总是等于加入到物体中的总能量这一点仍然正确。因而我们可以看到上面所推导的质量守恒等价于能量守恒。所以严格地说，像牛顿力学中那样的非弹性碰撞在相对论中是不存在的。按照牛顿力学两个物体碰撞后组成一个新的质量为$2m_0$的物体是毋庸置疑的，它与将这两个物体慢慢地放在一起所形成的物体毫无区别。当然，从能量守恒定律我们知道，在这个新的物体内有较多的动能，但是按照牛顿定律，这并不影响质量。然而现在我们看到这是不可能的；因为在碰撞中包含了动能，结果所产生的物体的质量将更大一些，因此，可以说，它是一个不同的物体。当我们轻轻地把物体组合在一起时，所产生的物体的质量为$2m_0$；当我们用力将物体组合在一起时，所产生的物体的质量更大一些。当质量不同时，我们就能识别出来。所以，在相对论中能量的守恒必定与动量守恒一同成立。

上面的讨论会产生一些有趣的结果。例如，假设我们有一个已测得质量为$m$的物体，如果发生某种情况使它分成两个相等的部分，各以速度$w$飞出，于是每个部分的质量将为$m_w$现在如果这两块裂片沿途碰上许多物体从而使速度变慢直至停止，那时它们的质量将为$m_0$。试问当它们停止时，给予其他物体多少能量？按照我们前面证明过的定理，每一部分提供的能量是$(m_w—m_0)c^2$。这么多能量就以热、势能或者别的什么形式留在其他物体里。由于$2m_w=M$，所以释放的能量是$E=(M—2m_0)c^2$。这个等式曾用来估计原子弹中的核裂变会释放多少能量（虽然许多碎片并不正好相等，而是近似相等）。铀原子的质量是已知的——在事前已测定过——铀原子裂变后产生的碘、等的质量也已知。这里的质量不是指原子运动时的质量，而是指它们静止的质量。换句话说，$M$和$m_0$都已知。这样，将有关两个数相减，我们就能计算如果$M$可以分裂为“两半”的话将会释放多少能量。由于这个理由，在所有的报纸上都曾将可怜的老爱因斯坦称为原子弹之“父”。当然，所有这些只是意味着，倘若我们告诉他会有什么反应发生的话，他就能在事先告诉我们会有多少能量被释放出来。一个铀原子裂变时应当释放的能量大约在第一次直接试验前六个月就估计出来了，后来能量真的释放时，立即有人直接作了测量（如果爱因斯坦公式不起作用，他们不管怎样也得测量能量），当他们测量出能量时，就不再需要这个公式了。当然，我们不应该贬低爱因斯坦，而是应该批评报界和许多报道文章对在物理和技术的发展过程中究竟是什么促成什么所提出的议论。至于怎样使事情更有效地和更迅速地出现，这完全是另一回事。

这个结果在化学上也同样有效，比方说，如果我们测定二氧化碳的质量，并与碳和氧的质量相比，我们就能算出，当碳和氧组成二氧化碳时，会释放出多少能量。这里唯一的麻烦在于质量上的差别是如此之小，以致这件事在技术上很难实现。

现在，我们回到这个问题上来：是否应当把$m_0c^2$加到动能上，从今以后就说一个物体的总能量是$m_0c^2$呢？首先，假如我们仍能在$M$内看出静止质量为$m_0$的组成部分，那么我们就可以说，复合物体的质量$M$中有一些是各组成部分的力学静止质量，有一些是各组成部分的动能，有一些是各组成部分的势能。但是我们发现，自然界中经历上述那种反应的各种粒子，在世界各国无论用什么方式来进行研究，都还没有能看到内部的组成成分。例如，K介子蜕变成两个$\pi$介子时，的确遵从定律式（\ref{Eq:I:16:11}），但是认为一个K介子由两个$\pi$介子组成则是无价值的观点，因为它也能蜕变为三个$\pi$介子！

因此我们有一种新的想法：我们毋须知道物体的内部结构；我们不能够，也不需要识别在粒子内部，哪一部分能量是粒子行将蜕变成的那些部分的静止能量。将一个物体的总能量$mc^2$分为内部组成部分的静止能量、动能和势能，既不方便，也常常不可能，我们只是简单地说粒子的总能量。我们对每个物体加上一个常数$m_0c^2$来“改变能量原点”，并且说，一个粒子的总能量是运动质量乘以$c^2$，而当物体静止时，其能量就是静止质量乘以$c^2$。

最后，我们发现速度$v$，动量$P$，总能量$E$能以一个相当简单的方式联系起来。很奇怪，以速率运动的物体的质量等于静止质量$m_0$除以$\sqrt{1-\frac{v^2}{c^2}}$这样一个公式却很少使用。相反，下面两个式子很容易证明，结果则很有用：
\begin{equation}
    \label{Eq:I:16:13}
    E^2-P^2c^2=m_0^2c^4
\end{equation}
和
\begin{equation}
    \label{Eq:I:16:14}
    Pc=\frac{Ev}{c}
\end{equation}